\section{Limitations}

Two kinds of assumption limit this study, and the second matters more than it may appear.

\begin{itemize}
\item \textbf{Parameters are estimates.} Every timing parameter (travel, door, boarding) and every behavioral
share (leaving spread, arrival distribution, transfer share, compliance, the two extra seats per car) is an
estimate, not a measurement. We varied several of them and the ordering of the measures held.
\item \textbf{Structure is assumed, and was not varied.} The simulator fixes how students behave in a queue,
how cars decide where to stop, how compliance is represented, and how demand is generated
(Table~\ref{tab:struct}). These choices are the same in every experiment, so no table in this paper measures their
influence, and they may shape the results as much as the parameters we vary. The robustness reported in
Section~\ref{sec:occ} should be read as robustness within one structure.
\item \textbf{No validation.} We calibrated $p_{occ}$ to match a qualitative description and compared nothing to
observed queues, wait times, or stop counts. Agreement with the description is not evidence that the structure is
right. The standard errors reflect replicate noise, not model error.
\item \textbf{Relative results only.} The model produces comparisons between measures that we believe are
reasonable, conditional on the points above. Absolute minutes and percentages should not be quoted as facts about
Mundelein Center.
\item \textbf{Unserved riders.} Earlier runs ended at 3{,}600~s and dropped riders still waiting, which understated
congestion at the highest occupancies. The cutoff is now a backstop that no scenario reaches, unserved riders would
count as late, and unserved counts are reported (Section~\ref{sec:future}, item~2). The accounting itself has not yet
been tested independently.
\item \textbf{Comparisons are not tested.} Measures are compared by their per-scenario means and standard errors,
not by confidence intervals on paired differences, and no outcome other than the ones listed in the Metrics
subsection is reported.
\item \textbf{One passing period.} We simulate one changeover from an empty building. Backlog from the previous
changeover would make the baseline worse.
\item \textbf{Uniform schedule.} Without a room-by-block schedule we assume every room is in use with probability
$p_{occ}$ in every block and cannot distinguish heavy blocks from light ones.
\item \textbf{Simple dispatching.} The real controller is proprietary and may behave differently. We deliberately
did not model a full-car bypass or down-peak mode, because these were out of scope.
\item \textbf{Static students.} Students do not adapt by leaving earlier, taking the stairs, or skipping the
line when they see it. This is one of the structural assumptions with the clearest real-world counterpart, since
adaptation would reduce both waits and lateness in the baseline.
\end{itemize}
